\szotaritem{Alagút }{  Olyan speciális úttípus, amelyben a hó és a jég nem hullik a csapadékmentesség miatt.
 }
\szotaritem{Autó (NPC) }{  A város lakóit szállító, önműködő jármű, amely a legrövidebb úton közlekedik a munkahely és a lakás között.
 }
\szotaritem{Baleseti ráta }{  A nehézségi foktól függő mutatószám, amely meghatározza a gépjárművek jeges úton történő összeütközésének valószínűségét.
 }
\szotaritem{Biokerozin }{  A sárkány fej működtetéséhez szükséges speciális üzemanyag, amely a telephelyen tölthető fel.
 }
\szotaritem{Buszvezető }{  A játékos egyik választható szerepköre, akinek célja a legtöbb kör megtétele a végállomások között.
 }
\szotaritem{Csomópont }{  A városi úthálózatot reprezentáló gráf azon pontjai, ahol az utak (élek) találkoznak, és ahol a forgalom várakozhat.
 }
\szotaritem{Egységnyi hó/jég }{  A csapadék belső mennyiségi egysége, amely a sűrűségtől függően eltérő fizikai vastagságot (centimétert) jelenthet az úton.
 }
\szotaritem{Feltört jég }{  A jégtörő fej által szétzúzott, de el nem takarított jégtörmelék, amely az autók áthaladása miatt újra veszélyes jégpáncéllá állhat össze.
 }
\szotaritem{Hányó fej }{  Olyan nagyteljesítményű takarítóeszköz, amely a havat és a feltört jeget az úttól messzire repíti.
 }
\szotaritem{Híd }{  Olyan úttípus, amelyen az alagúttal ellentétben folyamatosan és egyenletesen gyarapodik a hóréteg.
 }
\szotaritem{Hó }{  Az utakon folyamatosan és egyenletesen gyűlő csapadék, amely a gépjárművek súlya alatt jéggé tömörödik.
 }
\szotaritem{Hókotró }{  A takarító játékos által irányított jármű, amely különböző fejekkel felszerelve tisztítja az utakat.
 }
\szotaritem{Jégpáncél }{  Az autók és buszok áthaladása miatt tömörödött hóréteg, amely csúszásveszélyt okoz.
 }
\szotaritem{Jégtörő fej }{  A hókotróra szerelhető eszköz, amely a szilárd jégpáncélt darabokra töri, de nem távolítja el az útról.
 }
\szotaritem{Kijárási tilalom }{  A játék egyik végállapota (vereség), amely akkor következik be, ha a jeges utakon túl sok baleset történik a városban.
 }
\szotaritem{Körökre osztott működés }{  A játékmenet alapja, ahol a játékosok és az NPC-k egymás után, meghatározott sorrendben teszik meg lépéseiket.
 }
\szotaritem{Olvadási idő }{  A sózott útszakaszokon a hó és jég eltűnéséhez szükséges időtartam, amely egyenesen arányos a rétegvastagsággal.
 }
\szotaritem{Pass 'N Play }{  Olyan lokális többjátékos mód, amelyben a játékosok ugyanazon a számítógépen, egymást váltva hajtják végre a köreiket.
 }
\szotaritem{Sárkány fej }{  Gázturbinával felszerelt eszköz, amely biokerozin égetésével azonnal elolvasztja a havat és a jeget.
 }
\szotaritem{Sáv }{  Az utak belső felosztása; a hókotró egyszerre csak egy ilyen egységet tud megtisztítani, a forgalom pedig ezeken halad keresztül.
 }
\szotaritem{Só }{  A sószóró fej működéséhez szükséges fogyóeszköz, amely megakadályozza a hó megmaradását az úton.
 }
\szotaritem{Söprő fej }{  Olyan alapvető takarítóeszköz, amely a havat és a feltört jeget a menetirány szerinti jobbra szomszédos sávba tolja.
 }
\szotaritem{Sószóró fej }{  Olyan eszköz, amely só alkalmazásával olvasztja el a csapadékot és ideiglenesen tisztán tartja az útfelületet.
 }
\szotaritem{Súly (útszakaszé) }{  Az egyes útszakaszokhoz rendelt érték, amely meghatározza a takarításért kapható jutalom mértékét.
 }
\szotaritem{Takarító }{  A játékos azon szerepköre, amelyben a hókotró flottát és az üzemanyag-készleteket menedzseli.
 }
\szotaritem{Telephely }{  A városban található bázis, ahol a hókotrók felszereléseit cserélni, újakat vásárolni és üzemanyagot vételezni lehet.
 }
\szotaritem{Úthálózat }{  A várost reprezentáló gráf, ahol a csomópontok a kereszteződéseket, az élek pedig a különböző típusú utakat jelentik.
 }
\szotaritem{Végállomás }{  A buszok útvonalának két végpontja, amelyek között a buszvezetőknek a lehető legtöbb kört kell teljesíteniük. }
